% TOP_PROBLEM: {"id": 600010, "title": "Baker's Dozen — Unbounded unicycle tracks", "category_id": 6, "category": {"id": 6, "name": "geometry", "display_name": "Geometry", "description": "Euclidean and non-Euclidean geometry, geometric structures.", "slug": "geometry", "order_index": 6, "created_at": "2026-07-31T15:26:25.671Z"}, "status": "partially_solved", "problem_number": "AMR-005-0010", "set_id": 13}
% FIRST_POSTED: 2026-10-01
% ENTRY_KIND: statement_only
% UnsolvedMath ID 600010; source catalogue code AMR-005-0010
% !TeX program = lualatex
% Definition and sources only; this file is not an LLM solution attempt.
\documentclass[11pt,a4paper]{article}
\usepackage[margin=25mm]{geometry}
\usepackage{fontspec}
\setmainfont{lmroman10-regular.otf}[BoldFont=lmroman10-bold.otf,
 ItalicFont=lmroman10-italic.otf,BoldItalicFont=lmroman10-bolditalic.otf]
\usepackage{amsmath,amssymb,mathtools}
\usepackage{unicode-math}
\setmathfont{latinmodern-math.otf}
\AtBeginDocument{\let\setminus\smallsetminus}
\usepackage{xurl,hyperref}
\hypersetup{colorlinks=true,urlcolor=blue}
\setcounter{secnumdepth}{0}
\setlength{\parindent}{0pt}
\setlength{\parskip}{5pt}
\setlength{\emergencystretch}{3em}
\begin{document}
\section*{Baker's Dozen — Unbounded unicycle tracks}\label{600010}
{\small UnsolvedMath ID 600010 · AMR-005-0010 · Geometry · AMR Open Problem Lists}
\subsection{Definitions and mathematical statement}
Choose a smooth regular plane arc $\gamma_0:[0,L_0]\to\mathbb R^2$ parametrized
by arclength, with endpoints $(0,0)$ and $(1,0)$. At both endpoints its first
derivative is $(1,0)$ and all higher derivatives vanish. For any oriented
regular arc $\alpha$, define the tangent-offset arc
\[
 T(\alpha)(s)=\alpha(s)+\frac{\alpha'(s)}{\|\alpha'(s)\|}.
\]
Set $\gamma_{j+1}=T(\gamma_j)$, reparametrizing by arclength before the next
iteration. The endpoints of $\gamma_j$ are $(j,0)$ and $(j+1,0)$, and their
endpoint jets agree. Concatenating the arcs in this order gives a smooth
forward curve $\mathcal T$. The unit tangent, rather than an arbitrary-speed
derivative, is essential to this geometric construction.

Unless $\gamma_0$ is the horizontal straight segment, the source conjectures
all three conclusions: the heights on $\mathcal T$ are not bounded in any
horizontal strip; its image is not the graph of a single-valued function of
the horizontal coordinate; and its parametrization is not injective, so two
different traversal times give a self-intersection. The last condition is
stronger than simply ceasing to be a graph.
\subsection{Short English statement}
Repeatedly move an arc one unit along its tangent and join the resulting
arcs. Must every nonstraight seed eventually produce unbounded vertical
excursions, multiple heights at one horizontal position, and self-intersection?
\subsection{Sources}
\begin{itemize}
\item[S1] ULAM.AI and the UnsolvedMath Contributors, supplied problems.json, record 600010 (AMR-005-0010), AMR Open Problem Lists. Catalogue identity and source transcription; CC BY 4.0.
\url{https://huggingface.co/datasets/ulamai/UnsolvedMath}.
\item[S2] Tabachnikov - A Baker's Dozen of Problems (2015), Conjecture 3. Original source indexed by AMR-005-0010.
\url{https://amj.math.stonybrook.edu/html-articles/Files-2015-2024/14-01/}.
\end{itemize}
\end{document}
